%-----------------------------------------------------------------------
% This section is designed to be appended to bancroft_legal_tender.tex
% or included as a standalone document.
% It traces the Gold Clause Cases of 1935 back to the Legal Tender Cases
% discussed in the Bancroft document.
%-----------------------------------------------------------------------

\documentclass[12pt]{article}

\usepackage[margin=1.25in]{geometry}
\usepackage{setspace}
\usepackage{hyperref}
\usepackage{booktabs}
\usepackage{footnote}
\usepackage[T1]{fontenc}

\hypersetup{
    colorlinks=true,
    linkcolor=blue,
    citecolor=blue,
    urlcolor=blue
}

\doublespacing

\title{From the Legal Tender Cases to the Gold Clause Cases:\\
Bancroft's Nightmare Fulfilled, 1884--1935}
\author{Notes prepared from primary sources and Edwards (2018)}
\date{\today}

\begin{document}

\maketitle
\thispagestyle{empty}

\newpage
\tableofcontents
\newpage

%-----------------------------------------------------------------------
\section{Introduction: A Second Act for the Legal Tender Precedent}
%-----------------------------------------------------------------------

George Bancroft died in 1891, seven years after \textit{Juilliard v.\
Greenman} (1884) established that the legal tender power of Congress was
a permanent attribute of sovereignty, not limited to wartime emergencies.
He had warned that the decision ``retroactively undermined the
constitutional basis for the honor that Grant and the Republicans had
claimed in 1869'' when they pledged gold repayment of the 5-20 bonds.
But Bancroft could not have foreseen how completely the jurisprudential
machinery he denounced would be deployed a half-century later, when
Franklin Roosevelt abrogated all gold clauses in public and private
contracts, raised the price of gold from \$20.67 to \$35 per ounce, and
the Supreme Court was asked to decide whether the government could
repudiate its own explicit promises to pay in gold coin ``of the present
standard of value.''

The Gold Clause Cases of February~18, 1935---\textit{Norman v.\
Baltimore \& Ohio Railroad Co.}\ (294 U.S.\ 240), \textit{Nortz v.\
United States} (294 U.S.\ 317), and \textit{Perry v.\ United States}
(294 U.S.\ 330)---are the direct constitutional descendants of the
Legal Tender Cases that Bancroft excoriated. The majority opinions of
Chief Justice Hughes relied explicitly and repeatedly on \textit{Knox
v.\ Lee} and \textit{Juilliard v.\ Greenman} as foundational
precedents. The dissent of Justice McReynolds attempted to distinguish
the Legal Tender Cases from the Roosevelt gold program---and in doing
so essentially adopted the position that Bancroft, and Justice Field
before him, had articulated decades earlier. The episode is the subject
of Sebastian Edwards's \textit{American Default: The Untold Story of
FDR, the Supreme Court, and the Battle over Gold} (2018), which
provides a detailed economic and political narrative of the crisis.

%-----------------------------------------------------------------------
\section{The Roosevelt Gold Program, 1933--1934}
%-----------------------------------------------------------------------

The sequence of executive and legislative actions that culminated in
the Gold Clause Cases was breathtaking in its speed and scope:

\begin{itemize}

\item \textbf{March~6, 1933}: President Roosevelt declared a ``bank
holiday,'' citing heavy withdrawals of gold and speculative activity in
foreign exchange. The Secretary of the Treasury was instructed to make
payments in gold only under license.

\item \textbf{March~9, 1933}: The Emergency Banking Relief Act
confirmed the President's orders, amended the Trading with the Enemy
Act of 1917 to authorize regulation of gold transactions, and empowered
the Secretary of the Treasury to require surrender of all gold coin,
gold bullion, and gold certificates.

\item \textbf{April~5, 1933}: Executive Order~6102 required all persons
to deliver gold coin, bullion, and certificates to Federal Reserve
Banks by May~1, receiving ``an equivalent amount of any other form of
coin or currency'' in exchange.

\item \textbf{May~12, 1933}: Section~43 of the Agricultural Adjustment
Act authorized the President to reduce the weight of the gold dollar by
up to 50 percent and declared that all coins and currencies ``shall be
legal tender for all debts public and private.''

\item \textbf{June~5, 1933}: The Joint Resolution (Pub.\ Res.\ 73--10)
declared that gold clauses in all obligations, ``heretofore or
hereafter incurred,'' were ``against public policy'' and that all
obligations should be discharged ``upon payment, dollar for dollar, in
any coin or currency which at the time of payment is legal tender for
public and private debts.'' This applied explicitly to obligations
\textit{of} the United States as well as private contracts.

\item \textbf{January~30, 1934}: The Gold Reserve Act transferred
ownership of all monetary gold to the United States, provided that the
President should not fix the gold dollar at more than 60 percent of its
then-existing weight, and established a stabilization fund from the
resulting ``profit.''

\item \textbf{January~31, 1934}: By presidential proclamation, the
weight of the gold dollar was reduced from 25.8~grains to
15\,5/21~grains, nine-tenths fine---a devaluation of approximately
41~percent. The Treasury booked a paper profit of
\$2,800,000,000.

\end{itemize}

Edwards (2018) emphasizes that for months Roosevelt manipulated the
gold price through daily purchases, creating enormous uncertainty in
markets, and that the administration was prepared to defy the Supreme
Court if the Gold Clause Cases went against it. Roosevelt drafted
executive orders to close the stock exchanges and prepared a radio
address declaring that ``if the policy of the government \ldots\ is to be
irrevocably fixed by decisions of the Supreme Court,'' then ``the people
would have ceased to be their own rulers.'' Attorney General Homer
Cummings advised that the Court should be immediately packed to secure
a favorable ruling---foreshadowing the court-packing crisis of 1937.

%-----------------------------------------------------------------------
\section{The Three Cases}
%-----------------------------------------------------------------------

\subsection{\textit{Norman v.\ Baltimore \& Ohio Railroad Co.}\
(Private Gold Clauses)}

The lead case involved a \$22.50 coupon from a Baltimore \& Ohio Railroad
bond dated February~1, 1930, which promised payment ``in gold coin of
the United States of America of or equal to the standard of weight and
fineness existing on February~1, 1930.'' After devaluation, the gold
value of the coupon was \$38.10, but the railroad paid only \$22.50 in
legal tender currency. The bondholder sued for the difference.

Chief Justice Hughes, writing for a 5--4 majority (joined by Brandeis,
Stone, Roberts, and Cardozo), upheld the Joint Resolution as applied to
private contracts. The opinion's constitutional reasoning rested
squarely on the Legal Tender Cases.

\subsection{\textit{Nortz v.\ United States} (Gold Certificates)}

Nortz held \$106,300 in gold certificates---essentially warehouse
receipts for gold coin deposited with the Treasury---and was compelled
by Executive Order~6102 to surrender them for their face value in
currency. He sought an additional \$64,000 representing the loss from
devaluation. The Court held that since he could not lawfully have held
the gold itself, he had suffered no compensable damage.

\subsection{\textit{Perry v.\ United States} (Government Gold Bonds)}

Perry owned a \$10,000 Fourth Liberty Loan bond issued in 1918, which
promised payment of principal and interest ``in United States gold coin
of the present standard of value.'' After devaluation, he demanded
\$16,931.25 (at \$1.69 per dollar of the bond's face value). The
Treasury refused to pay more than \$10,000 in legal tender.

Hughes's opinion in \textit{Perry} reached a remarkable conclusion: the
Joint Resolution was \textit{unconstitutional} insofar as it purported
to abrogate the government's own gold clause obligations---Congress
could not use its power to regulate money so as to repudiate obligations
it had issued under the power to borrow ``on the credit of the United
States.'' But the bondholder nonetheless had \textit{no remedy}, because
he could not show actual damage: since gold could no longer be held or
exported, the dollars he received had the same purchasing power as the
gold dollars he was promised. The government's repudiation was declared
unconstitutional \textit{and} consequence-free.

Justice Stone concurred separately, agreeing with the result but
deploring the majority's willingness to pronounce on the
unconstitutionality of the Resolution when no damages could be
recovered: ``It will not benefit this plaintiff, to whom we deny any
remedy, to be assured that he has an inviolable right to performance
of the gold clause.''

%-----------------------------------------------------------------------
\section{The Legal Tender Cases as Foundation for the Gold Clause
Decisions}
%-----------------------------------------------------------------------

The connection between the 1870s Legal Tender Cases and the 1935 Gold
Clause Cases was not incidental---it was the explicit doctrinal
foundation on which the majority built its reasoning and against which
the dissent framed its protest.

\subsection{Hughes's Majority: \textit{Knox} and \textit{Juilliard} as
Controlling Precedent}

In \textit{Norman}, Hughes wrote that it was ``unnecessary to review
the historic controversy as to the extent of this power, or again to
go over the ground traversed by the Court in reaching the conclusion
that the Congress may make Treasury notes legal tender in payment of
debts previously contracted,'' citing \textit{Knox v.\ Lee} (12 Wall.\
457) and \textit{Juilliard v.\ Greenman} (110 U.S.\ 421) as
establishing the relevant ``settled principles.'' He then laid out
what he called the ``postulates'' of the Legal Tender decisions:

\begin{enumerate}

\item The source of congressional authority over currency lies not in
the coinage clause alone but in ``all the related powers conferred upon
the Congress and appropriate to achieve `the great objects for which the
government was framed'---`a national government, with sovereign
powers.'\,'' This was a direct quotation from \textit{Knox v.\ Lee} and
\textit{McCulloch v.\ Maryland}.

\item ``The broad and comprehensive national authority over the subjects
of revenue, finance, and currency is derived from the aggregate of the
powers granted to the Congress''---language taken directly from Justice
Gray's opinion in \textit{Juilliard}.

\item Congress has power ``to enact that the government's promises to
pay money shall be, for the time being, equivalent in value to the
representative of value determined by the coinage acts''---quoted from
\textit{Knox v.\ Lee} at page~553.

\item Congress may ``impress upon [its obligations] such qualities as
currency for the purchase of merchandise and the payment of debts, as
accord with the usage of sovereign governments''---quoted from
\textit{Juilliard} at page~447.

\item Contracts ``must be understood as having been made in reference to
the possible exercise of the rightful authority of the government, and
\ldots\ no obligation of a contract can extend to the defeat of that
authority''---from \textit{Knox v.\ Lee} at pages 549--551.

\item The Fifth Amendment does not prohibit indirect consequences of
legitimate monetary legislation, just as it does not prohibit the
effects of a new tariff or an embargo: ``whoever supposed that, because
of this, a tariff could not be changed?''---again directly from
\textit{Knox v.\ Lee}.

\end{enumerate}

Hughes also discussed the companion Legal Tender era cases---notably
\textit{Bronson v.\ Rodes} (7 Wall.\ 229) and \textit{Trebilcock v.\
Wilson} (12 Wall.\ 687)---which had upheld contracts for payment in
gold coin when two kinds of money were simultaneously in lawful
circulation. But Hughes turned these decisions against the gold clause
bondholders: those cases had arisen when ``gold was still in circulation
and no act of the Congress prohibiting the enforcement of such clauses
had been passed.'' The situation in 1935 was fundamentally different:
Congress had \textit{explicitly} prohibited gold transactions and had
established a single, uniform legal tender currency. Hughes noted that
Justice Bradley, in his concurrence in \textit{Knox v.\ Lee}, had
expressed the view that Congress \textit{could} have invalidated gold
clauses even in the 1860s and had merely chosen not to do so---a
``prophetic'' position now vindicated by the Joint Resolution.

In the \textit{Perry} opinion, Hughes similarly cited the Legal Tender
Cases for what they established and for how the present situation
differed: ``The case is not the same as if gold coin had remained in
circulation. That was the situation at the time of the decisions under
the legal tender acts of 1862 and 1863.'' The footnotes invoked Justice
Strong---who had written the majority in \textit{Knox v.\ Lee}---and
his subsequent opinion in the Sinking Fund Cases regarding the binding
character of the government's own contracts, as well as Alexander
Hamilton's 1795 statement that ``when a government enters into a
contract with an individual, it deposes, as to the matter of the
contract, its constitutional authority, and exchanges the character of
legislator for that of a moral agent.''

\subsection{McReynolds's Dissent: Distinguishing \textit{Knox} from
Roosevelt's Program}

Justice McReynolds's passionate dissent---delivered orally with such
vehemence that, according to contemporary accounts, he departed from
his text to exclaim ``this is Nero at his worst; the Constitution is
gone''---directly engaged the Legal Tender Cases and attempted to
distinguish them from Roosevelt's gold program. His argument
anticipated, and in some respects echoed, precisely the critique that
Bancroft had mounted against \textit{Juilliard} fifty years earlier.

McReynolds quoted the Legal Tender Cases at length, emphasizing the
passage from \textit{Knox v.\ Lee} that:
\begin{quote}
The legal tender acts do not attempt to make paper a standard of value.
We do not rest their validity upon the assertion that their emission is
coinage, or any regulation of the value of money; nor do we assert that
Congress may make anything which has no value money. What we do assert
is that Congress has power to enact that the government's promises to
pay money shall be, for the time being, equivalent in value to the
representative of value determined by the coinage acts, or to multiples
thereof.
\end{quote}

The critical words, for McReynolds, were ``for the time being.'' The
greenbacks issued during the Civil War were \textit{promises to pay
dollars}---``a pledge of the national credit''---with the expectation,
``ultimately realized,'' that ``in due time, they would be redeemed in
standard coin.'' And indeed they were, in 1879. ``The purpose was to
meet honorable obligations, not to repudiate them.'' The Legal Tender
Cases, McReynolds argued, provided ``no support to what has been
attempted here,'' because Roosevelt's program was fundamentally
different in character: ``The plan disclosed is \ldots\ primarily
designed to destroy private obligations, repudiate national debts, and
drive into the Treasury all gold within the country in exchange for
inconvertible promises to pay, of much less value.''

McReynolds also explicitly invoked the distinction that both the
Legal Tender dissents and Bancroft had drawn between emergency
wartime measures and permanent peacetime powers. He devoted a full
paragraph to showing that during the Civil War ``there was no coin---specie---in
general circulation in the United States between 1862 and 1879'' and
that ``both gold and silver were treated in business as commodities,''
so that the Legal Tender Cases ``arose during that period'' and could
not be mechanically extended. The key sentence was: ``This Court has
not heretofore ruled that Congress may require the holder of an
obligation to accept payment in subsequently devalued coins, or
promises by the government to pay in such coins.''

The Wikipedia entry on the Gold Clause Cases reports McReynolds's
position in summary form: ``McReynolds distinguished the cases at hand
from the Legal Tender Cases, arguing that in the earlier cases the
government sought to continue operating until it could meet its
obligations, while the Roosevelt administration apparently sought to
nullify them.'' This distinction---between using paper money as a
temporary bridge to future specie payment, and using paper money as an
instrument for permanent repudiation---was precisely the distinction
that Bancroft had made in his reading of the Legal Tender Cases and the
5-20 bond controversy.

%-----------------------------------------------------------------------
\section{Structural Parallels: From the 5-20s to the Liberty Bonds}
%-----------------------------------------------------------------------

The parallels between the 5-20 bond episode of the 1860s and the Gold
Clause Cases of the 1930s are striking and, in several respects,
structurally identical.

\subsection{The Ambiguity of the Promise}

In the 5-20 episode, the Act of February~25, 1862 specified that
\textit{interest} would be paid in coin but was silent on
\textit{principal}. Chase bridged the gap with reputational commitment,
and the ambiguity was exploited by Pendleton's ``Ohio Idea'' before
being resolved by the Public Credit Act of 1869.

In the Liberty Bond episode, the bonds explicitly promised payment of
\textit{both} principal and interest ``in United States gold coin of the
present standard of value.'' But the Joint Resolution of June~5, 1933
declared this explicit promise ``against public policy'' and provided
for discharge ``dollar for dollar'' in any legal tender. The 1930s
episode was thus, in one sense, more audacious than Pendleton's
proposal: where Pendleton had proposed to exploit a \textit{silence} in
the statute, Roosevelt abrogated an \textit{explicit} promise.

\subsection{The Quantum of the Discrimination}

The potential fiscal discrimination in 1868 was the difference between
paying the 5-20 principal in gold (or gold-equivalent currency, as
greenbacks had depreciated to roughly 72 cents on the dollar at their
nadir) and paying it in depreciated greenbacks. The actual greenback
discount at the time of the ``Ohio Idea'' debate was roughly 28~percent.

The fiscal discrimination in 1934 was determined by the devaluation
itself: the gold dollar was reduced from 25.8~grains to 15\,5/21
grains, a reduction of approximately 41~percent. Perry's \$10,000
Liberty Bond, on his calculation, entitled him to \$16,931.25 in the
devalued currency. The government's ``profit'' from the gold
revaluation was booked at \$2.8 billion. McReynolds pointed out with
acid precision that, if the power to devalue were unlimited, ``a gold
dollar containing one grain of gold may become the standard, all
contract rights fall, and huge profits appear on the Treasury
books\ldots\ perhaps \$20,000,000,000, maybe enough to cancel the
public debt, maybe more!''

\subsection{Resolution by Political Power}

Both episodes were resolved by the political branch that happened to
prevail. The 5-20 ambiguity was settled by Grant's election and the
Public Credit Act---in favor of the bondholders. The gold clause
ambiguity was settled by Roosevelt's congressional majorities, his
executive orders, and a narrow 5--4 Supreme Court decision---against
the bondholders. The Legal Tender Cases---specifically \textit{Knox v.\
Lee} and \textit{Juilliard}---supplied the constitutional architecture
for the 1935 result.

\subsection{Court-Packing: Threat and Precedent}

The parallel extends to the institutional machinery of constitutional
resolution. In 1870--71, Grant's appointment of Justices Strong and
Bradley to existing vacancies produced the reversal of \textit{Hepburn
v.\ Griswold} in \textit{Knox v.\ Lee}. Whether this constituted
deliberate court-packing remains debated, but Bancroft treated it as
such. In 1935, Attorney General Cummings openly recommended packing the
Court to secure a favorable decision in the Gold Clause Cases---and
when the Court ruled 5--4 for the government, the immediate pressure
was relieved, but the broader confrontation culminated in Roosevelt's
Judicial Procedures Reform Bill of 1937 (the ``court-packing plan'').

Edwards (2018) documents the administration's contingency planning in
detail, including Roosevelt's drafted radio address and executive
orders to close the exchanges. The threat to the Court's institutional
independence was, if anything, more explicit in 1935 than it had been
in 1870.

%-----------------------------------------------------------------------
\section{Bancroft's Ghost: \textit{Juilliard} Comes Home to Roost}
%-----------------------------------------------------------------------

The Gold Clause Cases of 1935 vindicated Bancroft's worst fears in
almost every particular. Consider the structure of the argument:

\begin{enumerate}

\item \textbf{Sovereignty as source of power}: Bancroft had condemned
Justice Gray's reasoning in \textit{Juilliard} that ``the power to make
the notes of the government a legal tender \ldots\ being one of the
powers belonging to sovereignty in other civilized nations, and not
expressly withheld from congress by the constitution'' justified the
legal tender power. Hughes in \textit{Norman} relied on precisely this
reasoning, deriving congressional authority from ``the aggregate of the
powers granted to the Congress'' rather than from any express grant, and
treating the power to regulate money as coextensive with national
sovereignty.

\item \textbf{Implied powers trumping explicit contracts}: Bancroft had
warned that \textit{Juilliard}'s logic would ``change the whole
character of our government'' from one of enumerated powers to one
without limits. In \textit{Norman}, Hughes held that ``contracts,
however express, cannot fetter the constitutional authority of the
Congress'' and that parties ``cannot remove their transactions from the
reach of dominant constitutional power by making contracts about them.''
This was the Legal Tender doctrine applied not merely to implicit money
contracts (as in \textit{Knox v.\ Lee}) but to \textit{explicit} gold
covenants---exactly the extension that Bancroft and Field had feared.

\item \textbf{No effective remedy}: Bancroft had observed that the
Constitution's nominal protections were meaningless if the Court would
not enforce them. In \textit{Perry}, Hughes declared the abrogation of
the government's gold clause unconstitutional---and then denied Perry
any recovery, on the ground that he could not show actual damage. The
government had made it illegal to hold gold, thereby preventing the
bondholder from establishing the loss that the unconstitutional action
had inflicted. As McReynolds put it: ``Having made payment in this
metal impossible, the government cannot defend by saying that, if the
obligation had been met, the creditor could not have retained the gold;
consequently he suffered no damage because of the nondelivery.''

\item \textbf{The Fourteenth Amendment's public debt clause}: Hughes
invoked Section~4 of the Fourteenth Amendment---``the validity of the
public debt of the United States, authorized by law, shall not be
questioned''---as ``confirmatory of a fundamental principle'' that
applied to the Liberty Bonds. But having invoked this sweeping
principle, he denied it any operative force. Bancroft would have
recognized the pattern: the Constitution says what it says, the Court
acknowledges what the Constitution says, and then the Court permits
Congress to do what the Constitution forbids.

\end{enumerate}

%-----------------------------------------------------------------------
\section{Edwards's \textit{American Default} and the Economics of
Abrogation}
%-----------------------------------------------------------------------

Sebastian Edwards's \textit{American Default} (Princeton University
Press, 2018) provides the most comprehensive modern account of the Gold
Clause episode. Edwards frames the abrogation of the gold clauses as a
sovereign default---the only peacetime default by the United States in
its history---and analyzes its economic consequences while also
providing a detailed narrative of the legal and political maneuvering.

Several of Edwards's contributions are particularly relevant to the
themes developed in this and the accompanying document on Bancroft:

\begin{itemize}

\item \textbf{The scale of the default}: Edwards estimates the total
face value of outstanding gold clause obligations at approximately
\$75~billion in early 1933, a figure consistent with the estimates
presented to the Court. The abrogation represented a transfer of wealth
from creditors to debtors of staggering proportions.

\item \textbf{The parallel to the 5-20 episode}: Edwards explicitly
draws the connection between the 1860s debate over paying the 5-20
bonds in greenbacks and the 1930s abrogation of gold clauses. In both
cases, the government's obligation was denominated in a specific medium
of payment, and the government sought either to substitute a cheaper
medium or to redefine the unit of account. The economic logic was
identical: fiscal pressure generated a temptation to inflate away (or
redefine away) the real burden of debt.

\item \textbf{The contingency planning}: Edwards documents in detail
Roosevelt's preparations for an adverse Supreme Court ruling, including
executive orders, the drafted radio address, and the policy of
deliberate market manipulation designed to create pressure for the
government's position. This planning---unknown at the time---reveals
that the administration viewed the Gold Clause Cases as an existential
political confrontation, not merely a legal dispute.

\item \textbf{The international dimension}: Edwards places the episode
in the context of global abandonment of the gold standard in the 1930s
and argues that the abrogation of gold clauses was part of a broader
international pattern. This is relevant to Bancroft's critique of
\textit{Juilliard}'s reasoning from what ``other civilized nations''
do: by 1935, other nations \textit{were} defaulting on gold
obligations, and the universality of the practice---far from being
shocking---was treated by the majority as evidence that the power was a
legitimate attribute of sovereignty.

\end{itemize}

%-----------------------------------------------------------------------
\section{The Arc from 1862 to 1935: Ambiguity, Commitment, and
Repudiation}
%-----------------------------------------------------------------------

Taken together, the 5-20 bond episode, the Legal Tender Cases, Bancroft's
\textit{Plea}, and the Gold Clause Cases describe a seventy-year arc in
American fiscal constitutionalism whose internal logic is remarkably
coherent.

\begin{enumerate}

\item \textbf{1862}: Congress creates deliberate ambiguity in the 5-20
bonds (silent on principal) and in the Legal Tender Act (silent on the
constitutional scope of the power). Chase bridges both ambiguities with
reputational commitment.

\item \textbf{1869--1871}: The political settlement resolves the 5-20
ambiguity in favor of bondholders (Public Credit Act) but resolves the
constitutional ambiguity in favor of governmental power (\textit{Knox
v.\ Lee}).

\item \textbf{1884}: \textit{Juilliard v.\ Greenman} extends the legal
tender power to peacetime, establishing that it is a permanent
attribute of sovereignty.

\item \textbf{1886}: Bancroft publishes his \textit{Plea}, warning that
the \textit{Juilliard} precedent has given Congress unlimited power
over the medium of payment and that no contractual promise---including
the government's own---is safe.

\item \textbf{1933--1935}: Roosevelt exercises exactly the power that
Bancroft feared, abrogating all gold clauses including those in
government bonds. The Supreme Court, by a 5--4 vote, sustains the
program, relying explicitly on \textit{Knox v.\ Lee} and
\textit{Juilliard} as controlling precedent.

\end{enumerate}

The Legal Tender Cases thus served as the constitutional transmission
mechanism between the Civil War fiscal emergency and the New Deal
monetary revolution. Without \textit{Knox v.\ Lee}'s reversal of
\textit{Hepburn}, and without \textit{Juilliard}'s extension of the
legal tender power to peacetime, the Gold Clause Cases would have
lacked their doctrinal foundation. Hughes could not have written that
it was ``unnecessary to review the historic controversy'' and could not
have treated the Legal Tender decisions as ``settled principles.'' The
resolution of the 1860s ambiguity---in favor of governmental power, and
against the hard-money constitutionalism of Chase's \textit{Hepburn}
opinion, of Field's dissent, and of Bancroft's \textit{Plea}---made the
1935 result, if not inevitable, at least doctrinally available.

Yet the Gold Clause Cases also reveal a distinction within the Legal
Tender tradition that complicates the story. In \textit{Perry}, Hughes
drew a ``clear distinction between the power of Congress to control or
interdict the contracts of private parties when they interfere with the
exercise of its constitutional authority, and a power in Congress to
alter or repudiate the substance of its own engagements.'' The
government's gold clause obligations were \textit{different} from
private obligations precisely because they were exercises of the
borrowing power, and ``the binding quality of the promise of the United
States is of the essence of the credit which is so pledged.'' This
distinction---between what the government can do to private contracts
and what it can do to its own---is the one principle from Bancroft's
tradition that survived into the 1935 majority opinion. But it survived
as a formal pronouncement without practical consequence: the government
was told it had acted unconstitutionally and was then told it owed
nothing.

The Perry paradox---unconstitutional but damage-free---is perhaps the
deepest illustration of the problem that Bancroft identified. A
constitutional commitment that the Court pronounces inviolable but
refuses to enforce is, in the language of modern commitment theory, a
soft constraint masquerading as a hard one. The bondholder learns that
his right is ``inviolable'' at exactly the moment when he learns that
it is worthless. Bancroft's fear that the Legal Tender decisions had
converted the Constitution's protections into ``vain promises'' found
its most exquisite confirmation in a decision that simultaneously
vindicated and nullified his central principle.

%-----------------------------------------------------------------------
\section*{References}
%-----------------------------------------------------------------------

\begin{description}
\setlength\itemsep{0.5em}

\item Bancroft, George. 1886. \textit{A Plea for the Constitution of
the United States of America Wounded in the House of Its Guardians}.
New York: Harper \& Brothers.

\item Edwards, Sebastian. 2018. \textit{American Default: The Untold
Story of FDR, the Supreme Court, and the Battle over Gold}. Princeton:
Princeton University Press.

\item Hall, George J. and Thomas J. Sargent. 2014. ``Fiscal
Discriminations in Three Wars.'' \textit{Journal of Monetary Economics}
61: 148--166.

\item Hughes, Charles Evans. 1935. Majority opinion in \textit{Norman
v.\ Baltimore \& Ohio Railroad Co.}, 294 U.S.\ 240 (February~18,
1935).

\item Hughes, Charles Evans. 1935. Majority opinion in \textit{Perry
v.\ United States}, 294 U.S.\ 330 (February~18, 1935).

\item McReynolds, James Clark. 1935. Dissenting opinion in \textit{Perry
v.\ United States}, 294 U.S.\ 330, at 361 (February~18, 1935).

\item Stone, Harlan Fiske. 1935. Concurring opinion in \textit{Perry
v.\ United States}, 294 U.S.\ 330, at 358 (February~18, 1935).

\item McKenna, Marian C. 2002. \textit{Franklin Roosevelt and the
Great Constitutional War: The Court-Packing Crisis of 1937}. New York:
Fordham University Press.

\item United States. Joint Resolution of June~5, 1933 (Pub.\ Res.\
73--10), 48 Stat.\ 112.

\item United States. Gold Reserve Act of January~30, 1934, 48 Stat.\
337.

\item United States. Emergency Banking Relief Act of March~9, 1933, 48
Stat.\ 1.

\item \textit{Hepburn v.\ Griswold}, 75 U.S.\ (8 Wall.) 603 (1870).

\item \textit{Knox v.\ Lee} and \textit{Parker v.\ Davis}, 79 U.S.\
(12 Wall.) 457 (1871).

\item \textit{Juilliard v.\ Greenman}, 110 U.S.\ 421 (1884).

\item \textit{Bronson v.\ Rodes}, 74 U.S.\ (7 Wall.) 229 (1869).

\item \textit{Trebilcock v.\ Wilson}, 79 U.S.\ (12 Wall.) 687 (1871).

\end{description}

\end{document}
